\subsection{Sionna Propagation Cache Mechanism}
\label{sec:sionna-cache-mechanism}

\subsubsection{Displacement-Based Cache Validity}
\label{subsec:cache-validity}

Performing the calculation using ray-tracing on Sionna RT is an computationally intensive calculation.
The goal of the cache mechanism is to reduce this computationally expensive calculation for every ns-3 propagation query.
When the receiver moved only a small distance from the previous position, the propagation information does not change significantly.
By using this idea, the cache validation is implemented based on spatial displacement rather than a fixed time interval.
In this section, how the cache is implemented is described in detail.

The cache stores the records of the each transmitter and receiver pair for ns-3 simulation.
Each cache entry is indexed by a canonical unordered node pair,
\begin{equation}
    \kappa(a,b)=\left(\min(I_a,I_b),\max(I_a,I_b)\right),
\end{equation}
where $I_a$ and $I_b$ are ns-3 node identifiers.
During the runtime, the cache stores the node positions $\mathbf{x}_a^0$ and $\mathbf{x}_b^0$, path loss, transmission delay, channel frequency reqponse (CFR), and the state for path whether it is in Line of Sight (LoS).

The distance between Tx and Rx is calculated using
\begin{equation}
    d_{a,b}^0=\left\|\mathbf{x}_a^0-\mathbf{x}_b^0\right\|_2 ,
\end{equation}
and the validity radius can be calculated using
\begin{equation}
    \Delta_{a,b}=
    \max\left(
    \Delta_{\min},
    \alpha \frac{0.886 d_{a,b}^0}{N_{\mathrm{tx,col}}}
    \right).
\end{equation}
where $\alpha$ is set to $0.4$ and $\Delta_{\min}$ is set to $0.1$ m as default values.

When the query for propagation delay and path loss happens for the Tx-Rx pair, the current posisitons  $\mathbf{x}_a(t)$ and $\mathbf{x}_b(t)$
are compared against the trace-time positions.
A valid cache entry is found if
\begin{equation}
    \left\|\mathbf{x}_a(t)-\mathbf{x}_a^0\right\|_2 < \Delta_{a,b}
    \quad\mathrm{and}\quad
    \left\|\mathbf{x}_b(t)-\mathbf{x}_b^0\right\|_2 < \Delta_{a,b}.
\end{equation}.
Ifboth inequalities are satisfied, the propagation delay and path loss and other data are returned directly.
Otherwise, the cache lookup is considered as a cache miss and the snapshot is refreshed.

\subsubsection{Snapshot Refresh}
\label{subsec:cache-refresh}

When the cache miss happen, it does not referesh only one link between Tx and Rx.
Instread, it refreshes the cache for all Tx-Rx pairs.
First, positions of every node which are using \texttt{SionnaMobilityModel} are synchronized into the Python Sionna scene.
Then it performs one batched calculation using ray-tracing for all Tx-Rx pairs by invoking \texttt{perform\_calculation} in \texttt{sionnart.py}.
This is designed in such a way to reduce overall computation using ray-tracing by calculating for all Tx-Rx pairs at once.
When the calculation are completes, it converts the returned records through pybind, and inserts them as new cache records in ns-3 propagation cache.
For each record, the cache stores
\begin{equation}
    \left(
    \tau_{a,b},\,
    L_{a,b},\,
    \mathbf{h}_{a,b},\,
    \mathbf{H}_{a,b},\,
    \mathbf{f},\,
    \mathrm{LoS}_{a,b},\,
    \mathbf{x}_a^0,\,
    \mathbf{x}_b^0,\,
    \Delta_{a,b}
    \right),
\end{equation}
where $\tau_{a,b}$ is delay, $L_{a,b}$ is path loss, $\mathbf{h}_{a,b}$ is the
scalar CFR, $\mathbf{H}_{a,b}$ is the MIMO CFR, and $\mathbf{f}$ is the
frequency vector.
The old cache records are removed only if the new results contains usable records.

\subsubsection{Lookup and Fallback Behavior}
\label{subsec:cache-lookup-fallback}

ns-3 query propagation related information for each Tx and Rx pair.
For a normal link, the cache lookup sequence is that it search the canonical key, test displacment validity, and return the first valid entry it found.
If no valid entry exists, the cache refreshes the full records and retris the same lookup sequence.

In Sionna RT, the link between nodes with same radio role, such as Tx--Tx or Rx--Rx, are not computable.
In such cases, the dummy entry with large path loss and delay are returned instead.
If Sionna RT does not return usable data for the requested link after refresh, the cache falls back to stochastic Friis model and constant constant-speed delay are returned for that lookup.

Moreover, before calling Sionna RT to refresh the cache, it checks estimated received power using Friis in order to avoid unnecessary computation expensive calculation.
If
\begin{equation}
    P_{\mathrm{Friis}}^{\mathrm{rx}} + M
    <
    P_{\mathrm{weak}},
\end{equation}
where $P_{\mathrm{weak}}=-95$ dBm and $M=6$ dB by default, the link is treated
as too weak to justify ray tracing.
In that case,  Friis loss and constant-speed delay are used directly.

\subsubsection{Derived MIMO Caches}
\label{subsec:derived-mimo-caches}

The cache records returned from Sionna side stores the normalized element-level MIMO CFR.
However, ns-3 needs derived quantities such as  port-level channel, matrices, PSD-scaled matrices, and effective gains after precoding.
These quantities depend not only on geometry but also on antenna objects, beamforming vectors, PSDs, and precoders.
Thus, these information are computed based on the returned element-level MIMO CFR and other information from the ns-3 side, and stored in the propagation entry as derived quantities.

The calculation can be done as followed.
For a given beam configuration, the element-level channel is projected onto
antenna ports.
For resource block $q$, transmit port $p_t$, and receive port
$p_r$, the port-level channel has the form
\begin{equation}
    G_{p_r,p_t}[q]
    =
    \sum_{i\in\mathcal{E}_{p_r}}
    \sum_{j\in\mathcal{E}_{p_t}}
    w_{r,i}\,w_{t,j}\,
    \widetilde{H}_{i,j}[q],
\end{equation}
where $\mathcal{E}_{p_r}$ and $\mathcal{E}_{p_t}$ are the receive and transmit
element sets associated with the ports, and $w_{r,i}$ and $w_{t,j}$ are the
beamforming weights.
This port CFR is cached using antenna identifiers,
beamforming-vector hashes, port counts, and link direction.

For spectrum propagation, the port CFR is scaled by the input PSD:
\begin{equation}
    C_{p_r,p_t}[q]
    =
    \sqrt{S[q]}\,G_{p_r,p_t}[q],
\end{equation}
where $S[q]$ is the PSD value on resource block $q$. The most recent
PSD-scaled matrix is cached using a PSD hash and PSD size.

When a precoding matrix $\mathbf{W}[q]$ is provided, the effective gain per
resource block is
\begin{equation}
    g[q]
    =
    \sum_{p_r}\sum_s
    \left|
    \sum_{p_t} G_{p_r,p_t}[q] W_{p_t,s}[q]
    \right|^2 .
\end{equation}
These gains are cached using the port-CFR key, precoder hash, precoder
dimensions, and output resource-block count.
This avoids repeated beamforming and precoding reductions while the underlying propagation entry remains valid.
